\section{Language model head} \label{sec:lm_head}
Additionally, we provide a description of the gradient computation for the loss function that involves the language model head. This computation entails performing a $\mathrm{softmax}$ operation over the entire vocabulary. If $\cV$ denotes the vocabulary set of the auxiliary model, and $\mE \in \RR^{\abs{\cV} \times \aux{D}}$ denotes the embedding matrix of the auxiliary model, we directly utilize the  embedding matrix for the auto-regressive loss in the \simulator. Additionally, we do not update the embedding matrix of the auxiliary model; instead, we solely backpropagate the gradients through the language model head. Recent work in~\citep{kumar2022finetune} has shown that keeping the embedding matrix fixed while updating the model can stabilize SGD. 
We demonstrate that the backpropagated gradients can be expressed as the combination of the language model head and a self-attention layer.



\begin{definition}[KL-loss gradient through auxiliary's language model head] \label{def:auxiliary_lm_back}
Given an embedding matrix $\mE \in \RR^{\abs{V} \times \aux{D}}$, the language model head takes in input $\vx \in \RR^{\aux{D}}$ and a target distribution $\vq \in \RR^{\abs{V}}$ and returns gradient $\losspartial{\vx} \in \RR^{\aux{D}}$, with $\losspartial{\vx} = \mE^{\top} \left( \mathrm{softmax} ( \mE \vx ) - \vq \right).$ 
\end{definition}


In the autoregressive loss on a sequence of tokens, the target output distribution at any position is the next occurring token. If  $\{\vx^{un}_t\}_{t=1}^{\aux{T}}$ denote the uncontextualized embeddings of a sequence of tokens after encoding them via the embedding matrix, and 
$\{\vx_t\}_{t=1}^{\aux{T}}$ denote their contextualized embeddings after passing through the auxiliary model, then the gradient $\losspartial{\vx_t}$ at any position $t$ can be simplified as $\mE^{\top}  \mathrm{softmax} ( \mE \vx_t ) - \vx^{un}_{t+1}.$ We illustrate the involved \simulator{} module w.r.t. an arbitrary position $t$.


\paragraph{\simulator{} autoregressive loss gradient module} 
The current embedding $\ve_t$ contains the contextualized embedding $\vx_t$ in its first $\aux{D}$ coordinates. Furthermore, $\ve_t$ includes the uncontextualized embedding $\vx^{un}_t$, copied from the input layer using residual connections. The prefix tokens ${ \vv_j }$ are assigned a value of $0$ and do not participate in the subsequent computations.

The loss computation can be decomposed into two sub-operations: (a) computing $\vy_t := \mE^{\top} \mathrm{softmax} ( \mE \vx_t )$, and (b) calculating $\losspartial{\vx_t} = \vy_t - \vx^{un}_{t+1}$.

For the first sub-operation, we use a feed-forward layer with $\mathrm{softmax}$ activation, with hidden and output weights $\mE$ and $\mE^{\top}$ respectively, that takes in the first $\aux{D}$ of $\ve_t$ and returns $\vy_t$ in the first $\aux{D}$ coordinates. We retain $\vx_t^{un}$ using a residual connection.

The final sub-operation can be interpreted as a \simulator{} self-attention layer. With $\ve_t$ containing both $\vy_t$ and $\vx_t^{un}$, we use a linear self-attention layer (\cref{def:self-attn_construct}) with two attention heads. The first attention head assigns an attention score of $1$ to pairs $\{(t, t+1)\}_{t \le \aux{T} - 1}$, while assigning an attention score of $0$ to the remaining pairs. At any position $t$, $-\vx^{un}_{t}$ is considered the value vector. The second attention head assigns an attention score of $1$ to pairs $\{(t, t)\}_{t \le \aux{T}}$, while assigning an attention score of $0$ to the remaining pairs. At any position $t$, $\vy_{t}$ is considered the value vector. The outputs of both attention heads are subsequently combined using a linear layer.


\begin{remark}
We conducted experiments using mean-squared loss and Quad loss \cite{saunshi2020mathematical}, which do not necessitate softmax computations for gradient computation. As an example, in the case of mean-squared loss, if our objective is to minimize $\frac{1}{2} \sum_{t=1}^{T} \norm{ \vx_t - \vx^{un}_{t+1} }^2$, the gradient can be computed as $\losspartial{\vx_t} = \vx_t - \vx^{un}_{t+1}$. Similarly, in the case of Quad loss, the gradient is $\losspartial{\vx_t} = \frac{1}{\abs{V}} \sum_{i} \ve_i - \vx^{un}_{t+1}$. However, in all of our language model experiments (\cref{sec:experiments}), both gradients resulted in minimal improvement in perplexity compared to the auxiliary model. Therefore, we continue utilizing the standard KL loss for optimization.
\end{remark}

\begin{remark}
For ease of implementation in the codebase, we utilize a dedicated loss module that takes in $\vy_t, \vx^{un}_{t+1}$ as input and directly computes $\losspartial{\vx_t} = \vy_t - \vx^{un}_{t+1}$.
\end{remark}








